\scanpage{73}
considered as a subgroup of the Brauer group consisting of the elements annihilated by 2. Thus \uline{the structure of} $Q(A)$ \uline{is known when one knows the groups} $H^{\handwritten{i}}(A,\underline{Z}/2\underline{Z})$ \uline{for} $i=1,2$, \uline{namely} $D$ \uline{and} $B$, \uline{and the cup-product} $D\times D\handwritten{\to}B$. The preceding statement shows \deleted{indeed}\inserted{\handwritten{even}} that the kernel of the surjective homomorphism (\handwritten{*}) is also the subgroup generated by the elements
\[
e_{a+u+v}-e_{a+u}-e_{a+v}+e_a
\]
for the elements $a,u,v$, in $D$ such that $uv=0$.

One concludes from this in particular that the associated graded of $Q(A)$ is isomorphic to the quotient of the symmetric algebra of $D$ over $\underline{Z}$, by \deleted{\unsure{an}} the ideal generated by the elements of order 2 of the form $ab-a'b'$, where $a,b,a',b'$, are as above, or what amounts to the, same, generated by the elements $uv$ of $\mathrm{Symm}^{\handwritten{2}}(D)$ for $u,v\quad D$ such that $u.v=0$ in $B^{\deleted{\unsure{x}}}$. (N.B. it would be necessary to clarify whether $\mathrm{grad}^{\handwritten{2}}Q(A)$ \deleted{is} $\handwritten{\to}B$ is \deleted{\unsure{very}}\handwritten{always} injective, in other words whether the kernel of $\mathrm{Symm}^{\handwritten{2}}(D)\quad B$ is generated by elements of the form \deleted{uv} $uv$).

One concludes from the \deleted{\unsure{whole|our}} structure theorem that if ${\unsure{\mathbb{H}|H}}^{\handwritten{\bullet}}$ is a graded ring in positive degrees, then giving an \deleted{additive}\inserted{\handwritten{augmented ring}} homomorphism from $Q(A)$ to $\handwritten{\underline{Z}\times(1+}\overset{\handwritten{\scriptscriptstyle\wedge}}{\mathbb{H}}^{\handwritten{+}}\handwritten{)}$ is equivalent to giving a homomorphism
\[
w^{\handwritten{1}}:D\handwritten{\longrightarrow}\mathbb{H}^{\handwritten{1}}
\]
such that $W^{\handwritten{1}}(a)W^{\handwritten{1}}(b)=W^{\handwritten{1}}(a')W^{\handwritten{1}}(b')$ when $a,b,a',b'$ are as above, or what amounts to the same, such that $W^{\handwritten{1}}(u)W^{\handwritten{1}}(v)=0$ if $u.v=0$ in $B$; \deleted{one} moreover $\mathrm{grad}\ Q(A)$ plays a universal role for its properties, i.e. giving $W$ is also equivalent to giving a \deleted{addi} homomorphism of graded rings
\[
\mathrm{grad}\ Q(A)\handwritten{\to}H^{\handwritten{\bullet}}\ .
\]
The most natural choice for $H^{\handwritten{\bullet}}$, from the ``topological'' point of view, seems
\[
H_{\deleted{\unsure{i|1}}}^{\handwritten{\bullet}}=H^{\handwritten{\bullet}}(A,\underline{Z}/2\underline{Z})\ ,
\]
equipped with the cup-product, which gives us \uline{Stiefel-Whitney} classes
\[
W^{\handwritten{i}}:Q(A)\handwritten{\to}H^{\handwritten{1}}(A,\underline{Z}/2\underline{Z})\ .
\]
Serre suggests \deleted{taking} also considering a finer definition of the Stiefel-Whitney classes, by replacing \deleted{\unsure{the grou|the group}} the maximal separable algebraic extension of $A$ by the subextension union of the finite subextensions of rank a power of 2; this would already give a ``finer'' $H^{\handwritten{2}}$, i.e. mapping to the usual $H^{\handwritten{2}}$ ...
% archive sheet number: 72
